\documentclass[10pt,a4paper]{amsart}
\usepackage[margin=19mm]{geometry}
\usepackage{amsmath,amssymb,mathtools}
\usepackage[scr=rsfs,cal=euler]{mathalfa}
\usepackage{xfrac}
\usepackage[unicode,hidelinks,bookmarks=false]{hyperref}
\newcommand{\ie}{i.e.\ }
\newcommand{\cf}{cf.\ }
\newcommand{\resp}{respectively\ }
\newcommand{\Subsection}{\subsection}
\providecommand{\nobreakdash}{\nobreak\hbox{-}}
\setlength{\emergencystretch}{2em}
\multlinegap0pt
\usepackage{bm}
\usepackage{bbm}
\usepackage{dsfont}
\usepackage{xspace}
\usepackage{cancel}
\usepackage{mathtools}
\usepackage{cleveref}
\usepackage{amsthm}
\makeatletter
\@ifundefined{thm}{\newtheorem{thm}{Thm}}{}
\@ifundefined{lem}{\newtheorem{lem}[thm]{Lemma}}{}
\@ifundefined{prop}{\newtheorem{prop}[thm]{Proposition}}{}
\makeatother
\newcommand{\Rt}{\mathsf{R}}                 % the D4 root system
\title[The Sphere Packing Problem in Dimension 4 and the Twenty-Fou]{The Sphere Packing Problem in Dimension 4 and the Twenty-Four-Cell Conjecture}
\author{Deep Bhattacharjee, Ushashi Bhattacharya, Priyabrata Mandal,, Shounak Bhattacharya}
\date{Source: arXiv:2609.25120v2 (CC BY 4.0)}
\begin{document}
\maketitle

\noindent\emph{Source and licence.} The Sphere Packing Problem in Dimension 4 and the Twenty-Four-Cell Conjecture. Version arXiv:2609.25120v2 (CC BY 4.0), distributed under the Creative Commons Attribution 4.0 International licence, as recorded in the arXiv licence metadata for this exact version and shown on the abstract page https://arxiv.org/abs/2609.25120v2. Full rights notice: licenses/arxiv-2609.25120-CC-BY-4.0.txt.\par\smallskip

\noindent\textit{Editorial scope.} Counted unit: the Prop whose source label is \texttt{prop:meet21} in the article (source0.tex lines 4152--4155, source label \texttt{prop:meet21}), delivered together with the article's own complete original proof of it (source0.tex lines 4157--4243, one \texttt{proof} environment of 87 lines). No mathematics is written, repaired or re-derived: statement and proof are the article's own text line for line. Required context retained uncounted: eq:couples (lines 4096--4102); lem:filled-support (lines 4106--4109); lem:filled-couple (lines 4118--4121). Every definition, display and label of those lines is retained unchanged. Omitted from this package and not redistributed: 1--4095, 4103--4105, 4110--4117, 4122--4151, 4156--4156, 4244--10670. Nothing inside a retained excerpt is omitted or altered: every retained body is delivered exactly as the article's lines are written, so no omission receipt of any kind accompanies this package. Citations: the retained excerpts carry no citation command at all, so no bibliography entry of the article is transcribed and no reference is redistributed. Reference adaptation: none was needed and none was made; every reference inside the delivered unit resolves inside it, so no unresolved reference or citation is produced. Printed slips kept literal: no printed word, symbol, hypothesis or number of the counted statement or of its proof was altered. Layout and class: the article's own class is \texttt{amsart} with packages including fontenc, lmodern, amsmath, amssymb, amsthm, amsfonts, mathtools, mathrsfs, enumitem, xcolor, booktabs, multirow, longtable, graphicx; the wrapper is the lane's pinned \texttt{amsart} editorial wrapper at 10pt on A4, which restates the theorem-like environments the retained text needs and adds the article's own macro definitions that the retained text needs (\texttt{\textbackslash Rt}), with extra wrapper packages (bm, bbm, dsfont, xspace, cancel, mathtools, cleveref). The delivered numbering is the unit's own. Assets: no retained span contains a figure environment, a graphics-inclusion command, a PDF-page inclusion, a table or a TikZ picture; no image file is copied into this package, the package declares no asset dependency and no image file is redistributed with it. Licence: version arXiv:2609.25120v2 of this article is distributed under the Creative Commons Attribution 4.0 International licence, as recorded in the arXiv licence metadata for this exact version and shown on the abstract page https://arxiv.org/abs/2609.25120v2. A full rights notice is delivered at licenses/arxiv-2609.25120-CC-BY-4.0.txt. Characters: every retained excerpt is pure ASCII, so the delivered engine, XeLaTeX with its default Latin Modern text font, reports no missing character.
\begin{equation}\label{eq:couples}
  \{1,2\}\,\big|\,\{3,4\},
  \qquad
  \{1,3\}\,\big|\,\{2,4\},
  \qquad
  \{1,4\}\,\big|\,\{2,3\},
\end{equation}

\begin{lem}[A filled support]\label{lem:filled-support}
If $\{i,j\}$ is filled for $S$, then every $v\in A(S)$ satisfies
$|x_i|+|x_j|\le t$, and hence $x_i^2+x_j^2\le\tfrac12$.
\end{lem}

\begin{lem}[A filled couple]\label{lem:filled-couple}
If both supports of one of the couples \eqref{eq:couples} are filled for
$S$, then $A(S)=\Rt\setminus S$.
\end{lem}

\begin{prop}[Twenty-one roots]\label{prop:meet21}
If a contact configuration $W$ satisfies $|W\cap\Rt|=21$ for some root
system $\Rt$, then $W\subseteq\Rt$ or $|W|\le22$.
\end{prop}

\begin{proof}
Let $S=W\cap\Rt$ and let $\rho_1,\rho_2,\rho_3$ be the missing roots. If
their supports occupy at most two of the six index pairs, some couple is
filled, so $A(S)=\{\rho_1,\rho_2,\rho_3\}$ by \cref{lem:filled-couple}
and $W\subseteq\Rt$. So assume the three supports are distinct and meet
each couple of \eqref{eq:couples} once. Choosing one index pair from each
couple gives eight possibilities, and inspection of the list shows that
four of them are \emph{stars}, three pairs sharing an index, as with
$\{1,2\},\{1,3\},\{1,4\}$, and four are \emph{triangles}, the three pairs
inside a triple of indices, as with $\{1,2\},\{1,3\},\{2,3\}$. Permuting
the coordinates reduces to those two.

Let $v\in A(S)$ and suppose $v$ is not a root. For each $m$ the set
$S\cup\{\rho_m\}$ misses only two roots, so $A(S\cup\{\rho_m\})$ consists
of roots by \cref{lem:filled-couple}, and therefore does not contain $v$.
Since $v\in A(S)$, the only condition that can fail at $S\cup\{\rho_m\}$
is the one at $\rho_m$ itself, so
\begin{equation}\label{eq:inside-sixty}
  \langle v,\rho_m\rangle>\tfrac12,\qquad m=1,2,3 .
\end{equation}

Suppose first that the supports form the star $\{1,2\},\{1,3\},\{1,4\}$.
Sign changes in the coordinates $2,3,4$ are symmetries of $\Rt$, so we
may take $\rho_m=(\varepsilon_me_1+e_m)/\sqrt2$ with
$\varepsilon_m\in\{\pm1\}$ for $m=2,3,4$. The supports $\{2,3\}$,
$\{2,4\}$ and $\{3,4\}$ are filled, so $|x_m|\le t$ for $m=2,3,4$ by
\cref{lem:filled-support}, and \eqref{eq:inside-sixty} reads
$\varepsilon_mx_1+x_m>t$. If $\varepsilon_m=1$ this gives
$x_1>t-x_m\ge0$, and if $\varepsilon_m=-1$ it gives $x_1<x_m-t\le0$; the
two cannot both occur, so the three signs agree, and after a sign change
in the first coordinate $\rho_m=(e_1+e_m)/\sqrt2$ for $m=2,3,4$. Those
three roots are pairwise at $60^\circ$.

Suppose instead that the supports form the triangle
$\{1,2\},\{1,3\},\{2,3\}$, and write
$\rho_1=(ae_1+be_2)/\sqrt2$, $\rho_2=(ce_1+de_3)/\sqrt2$,
$\rho_3=(fe_2+ge_3)/\sqrt2$ with signs $a,\dots,g$. The supports
$\{1,4\},\{2,4\},\{3,4\}$ are filled, so $|x_m|\le t$ for $m=1,2,3$.
Adding the three inequalities \eqref{eq:inside-sixty} gives
\[
  (a+c)\,x_1+(b+f)\,x_2+(d+g)\,x_3>3t ,
\]
and the left side is at most $2tk$, where $k$ counts the nonzero terms
among $a+c$, $b+f$, $d+g$. Hence $k=3$, that is $a=c$, $b=f$ and $d=g$,
and $\rho_1,\rho_2,\rho_3$ are again pairwise at $60^\circ$. Changing the
signs of the coordinates $1,2,3$ as needed makes the triple
$\{(e_1+e_2),(e_1+e_3),(e_2+e_3)\}/\sqrt2$, and the orthogonal map with
matrix $\tfrac12\bigl(\begin{smallmatrix}1&1&1&1\\1&1&-1&-1\\
1&-1&1&-1\\1&-1&-1&1\end{smallmatrix}\bigr)$, which permutes $\Rt$,
carries it to $\{(e_1+e_2),(e_1+e_3),(e_1-e_4)\}/\sqrt2$; one more sign
change in the fourth coordinate returns the star.

So in every case we may take $\rho_m=(e_1+e_m)/\sqrt2$ for $m=2,3,4$, and
it remains to read off $A(S)$ for that $S$. The supports $\{2,3\}$,
$\{2,4\}$, $\{3,4\}$ are filled, so $|x_2|+|x_3|$, $|x_2|+|x_4|$ and
$|x_3|+|x_4|$ are all at most $t$. The maximum of $p^2+q^2+r^2$ over
$\{p,q,r\ge0:\ p+q\le t,\ p+r\le t,\ q+r\le t\}$ is attained at a vertex
of that polytope, and its vertices are the origin, the three points
$te_k$ and $\tfrac t2(1,1,1)$, with values $0$, $t^2$ and $\tfrac34t^2$.
Hence
\begin{equation}\label{eq:doorway-coords}
  x_2^2+x_3^2+x_4^2\le t^2=\tfrac12,
  \qquad\text{so}\qquad
  x_1^2\ge\tfrac12 .
\end{equation}
Each support $\{1,m\}$ keeps three of its four roots, all but
$(e_1+e_m)/\sqrt2$, and these give $x_1-x_m\le t$ and $|x_m|-x_1\le t$.
If $x_1\le-t$ the second forces $x_m=0$ for $m=2,3,4$, hence $v=-e_1$,
which the first contradicts at $m=2$. So $x_1\ge t$, and then
$x_m\ge x_1-t\ge0$ for $m=2,3,4$: on $A(S)$ the first coordinate is at
least $t$ and the other three are nonnegative. Consequently, for
$v,v'\in A(S)$,
\begin{equation}\label{eq:doorway-pair}
  \langle v,v'\rangle=x_1x_1'+x_2x_2'+x_3x_3'+x_4x_4'
  \;\ge\;x_1x_1'\;\ge\;t^2=\tfrac12 ,
\end{equation}
with equality only if $x_1=x_1'=t$. That makes
\eqref{eq:doorway-coords} an equality, so $(x_2,x_3,x_4)$ is one of the
vertices $te_k$, and $v\in\{\rho_1,\rho_2,\rho_3\}$; the same holds for
$v'$.

Now $W\setminus S$ lies in $A(S)$ and its elements are pairwise at
$60^\circ$ or more. By \eqref{eq:doorway-pair} any two of them are
missing roots, so either every element of $W\setminus S$ is a root, in
which case $W\subseteq\Rt$, or $W\setminus S$ has a single element and
$|W|=22$.
\end{proof}

\end{document}
